\setcounter{section}{26}

  \zwischenueberschrift{Übungsaufgaben}

\inputaufgabe
{}
{

Die reelle Ebene sei mit der
\definitionsverweis {euklidischen Struktur}{}{}
und dem zugehörigen
\definitionsverweis {Levi-Civita-Zusammenhang}{}{}
versehen. Berechne

\mavergleichskette
{\vergleichskette
{ \nabla_FG
}
{ = }{
}
{  }{
}
{    }{
}
{   }{
}
}
{}{}{}
zu den Vektorfeldern

\mavergleichskettedisp
{\vergleichskette
{ F(x,y)
}
{ =} { \left(  x^2-y^3  , \,   xy+y^2                   \right)
}
{ } {
}
{ } {
}
{ } {
}
}
{}{}{}
und

\mavergleichskettedisp
{\vergleichskette
{ G(x,y)
}
{ =} { \left( e^{xy^2}   , \,   x^3-x^2y^5                   \right)
}
{ } {
}
{ } {
}
{ } {
}
}
{}{}{.}

}
{} {}

\inputaufgabe
{}
{

Es sei

\mavergleichskette
{\vergleichskette
{ I
}
{ \subseteq }{ \R
}
{  }{
}
{    }{
}
{   }{
}
}
{}{}{}
ein
\definitionsverweis {reelles Intervall}{}{}
und sei

\mavergleichskettedisp
{\vergleichskette
{ g(t)
}
{ =} { 3+t^2
}
{ } {
}
{ } {
}
{ } {
}
}
{}{}{,}
die wir als eine
\definitionsverweis {riemannsche Metrik}{}{}
auf $I$ interpretieren. Bestimme die
\definitionsverweis {vertikale Ableitung}{}{}
\zusatzklammer {zu dem zugehörigen
\definitionsverweis {Levi-Civita-Zusammenhang}{}{}} {} {}

\mathl{\nabla_\partial  \left(  t^3-  \sin  \left( t^2 \right)   \right)}{.}

}
{} {}

\inputaufgabegibtloesung
{}
{

Es sei

\mavergleichskette
{\vergleichskette
{ I
}
{ \subseteq }{ \R
}
{  }{
}
{    }{
}
{   }{
}
}
{}{}{}
ein
\definitionsverweis {reelles Intervall}{}{}
und sei
\maabbdisp {g} { I } {\R_+
} {}
eine positive differenzierbare Funktion, die wir als eine
\definitionsverweis {riemannsche Metrik}{}{}
auf $I$ interpretieren. Bestimme die
\definitionsverweis {horizontalen Schnitte}{}{}
zu dem zugehörigen
\definitionsverweis {Levi-Civita-Zusammenhang}{}{}
auf $I \times \R$
\zusatzklammer {siehe
[[Intervall/Riemannsche Struktur/Levi-Civita-Zusammenhang/Beispiel|Beispiel 26.4]]} {} {.}

}
{} {}

\inputaufgabe
{}
{

Es sei $M$ eine
\definitionsverweis {differenzierbare Mannigfaltigkeit}{}{}
derart, dass das
\definitionsverweis {Tangentialbündel}{}{}
$TM$ trivial ist, und sei auf $M$ über eine differenzierbare Trivialisierung

\mavergleichskette
{\vergleichskette
{ TM
}
{ \cong }{ M \times \R^n
}
{  }{
}
{    }{
}
{   }{
}
}
{}{}{}
\zusatzklammer {mithilfe des
\definitionsverweis {Standardskalarproduktes}{}{}
auf dem $\R^n$} {} {}
eine
\definitionsverweis {riemannsche Struktur}{}{}
gegeben
\zusatzklammer {siehe
[[Mannigfaltigkeit/Tangentialbündel trivial/Riemannsche Struktur/Aufgabe|Aufgabe 16.2]]} {} {.}
Zeige, dass der
\definitionsverweis {Levi-Civita-Zusammenhang}{}{}
auf $M$ der
\definitionsverweis {triviale Zusammenhang}{}{}
ist.

}
{} {}

\inputaufgabe
{}
{

Es sei ${\mathbb H}$ die
\definitionsverweis {Halbebene von Poincaré}{}{}
mit der zugehörigen
\definitionsverweis {riemannschen Struktur}{}{}
und dem zugehörigen
\definitionsverweis {Levi-Civita-Zusammenhang}{}{.}
Bestimme die
\definitionsverweis {vertikale Ableitung}{}{}
\mathl{\nabla_{\partial_1 + \partial_2 } (G)}{} zum Vektorfeld

\mavergleichskettedisp
{\vergleichskette
{ G(x,y)
}
{ =} { \left(  x^5-y^3+2xy  , \,   4 x^2+ y^3                   \right)
}
{ } {
}
{ } {
}
{ } {
}
}
{}{}{.}

}
{} {}

\inputaufgabegibtloesung
{}
{

Es sei ein
\definitionsverweis {linearer Zusammenhang}{}{}
$\nabla$ auf dem
\definitionsverweis {Tangentialbündel}{}{}
$TX$ einer
\definitionsverweis {differenzierbaren Mannigfaltigkeit}{}{}
$X$ gegeben. Es gebe
\definitionsverweis {Basisschnitte}{}{}

\mathl{V_1 , \ldots , V_n}{} mit

\mavergleichskettedisp
{\vergleichskette
{ \nabla_{V_i } V_j -  \nabla_{V_j } V_i
}
{ =} { [V_i,V_j]
}
{ } {
}
{ } {
}
{ } {
}
}
{}{}{}
für alle $i,j$. Zeige, dass $\nabla$
\definitionsverweis {torsionsfrei}{}{}
ist.

}
{} {}

\inputaufgabe
{}
{

Zeige, dass ein
\definitionsverweis {linearer Zusammenhang}{}{}
auf dem
\definitionsverweis {Tangentialbündel}{}{}
eines offenen Intervalls

\mavergleichskette
{\vergleichskette
{ I
}
{ \subseteq }{ \R
}
{  }{
}
{    }{
}
{   }{
}
}
{}{}{}
\definitionsverweis {torsionsfrei}{}{}
ist.

}
{} {}

\inputaufgabe
{}
{

Es sei

\mavergleichskette
{\vergleichskette
{ U
}
{ \subseteq }{ \R^n
}
{  }{
}
{    }{
}
{   }{
}
}
{}{}{}
\definitionsverweis {offen}{}{}
und sei darauf eine
\definitionsverweis {riemannsche Struktur}{}{}
durch die Funktionen $\left(  g_{ij}  \right)$ gegeben mit der
\definitionsverweis {inversen Matrix}{}{}

\mathl{\left(  g^{k l}  \right)}{.} Es sei $\nabla$ der durch die
\definitionsverweis {Christoffelsymbole}{}{}

\mavergleichskettedisp
{\vergleichskette
{ \Gamma_{ij}^k
}
{ =} { { \frac{   1  }{  2 } } \sum_{l = 1}^n g^{kl}(\partial_ig_{jl}+\partial_jg_{il}-\partial_lg_{ij})
}
{ } {
}
{ } {
}
{ } {
}
}
{}{}{}
gegebene
\definitionsverweis {lineare Zusammenhang}{}{,}
der durch

\mavergleichskettedisp
{\vergleichskette
{ \nabla_{\partial_i } \partial_j
}
{ =} { \sum_{k = 1}^n\Gamma_{ij}^k\,\partial_k
}
{ } {
}
{ } {
}
{ } {
}
}
{}{}{}
gekennzeichnet ist. Zeige in zwei Schritten, dass $\nabla$ die Koszul-Formel aus
[[Riemannsche Mannigfaltigkeit/Tangentialbündel/Linearer Zusammenhang/Metrisch und torsionsfrei/Koszul-Eigenschaft/Eindeutigkeit/Fakt|Lemma 26.10]]
erfüllt.
\aufzaehlungzwei {Für die Basisfelder $\partial_i$.
} {Für beliebige Vektorfelder $V,W,Z$.
}

}
{} {}

\inputaufgabegibtloesung
{}
{

Es sei $M$ eine
\definitionsverweis {riemannsche Mannigfaltigkeit}{}{}
und sei $\nabla$ ein
\definitionsverweis {linearer Zusammenhang}{}{}
auf $TM$. Es sei vorausgesetzt, dass $\nabla$ lokal für Basisfelder $\partial_1  , \ldots , \partial_n$ die Eigenschaft

\mavergleichskettedisp
{\vergleichskette
{ \partial_i \left\langle  \partial_j  ,  \partial_k   \right\rangle
}
{ =} { \left\langle \nabla_{\partial_i } \partial_j   ,  \partial_k   \right\rangle  + \left\langle  \partial_j  ,  \nabla_{\partial_i } \partial_k    \right\rangle
}
{ } {
}
{ } {
}
{ } {
}
}
{}{}{}
für alle $i,j,k$ erfüllt. Zeige, dass $\nabla$
\definitionsverweis {metrisch}{}{}
ist.

}
{} {}

Das Konzept eines metrischen Zusammenhangs gibt es nicht nur auf dem Tangentialbündel einer riemannschen Mannigfaltigkeit, sondern allgemeiner auf einem riemannschen Vektorbündel.

Es sei
\maabb {p} { E } { M
} {}
ein
\definitionsverweis {riemannsches Vektorbündel}{}{}
über einer
\definitionsverweis {differenzierbaren Mannigfaltigkeit}{}{}
$M$. Ein
\definitionsverweis {linearer Zusammenhang}{}{}
auf $E$ heißt
\definitionswort {metrisch}{,}
wenn auf jeder offenen Menge

\mavergleichskettedisp
{\vergleichskette
{ D_V \left\langle  s  , t  \right\rangle
}
{ =} { \left\langle \nabla_V s , t  \right\rangle  + \left\langle  s  ,  \nabla_Vt  \right\rangle
}
{ } {
}
{ } {
}
{ } {
}
}
{}{}{}
für ein
\definitionsverweis {stetig differenzierbares}{}{}
\definitionsverweis {Vektorfeld}{}{}
$V$ und stetig differenzierbare Schnitte $s,t$ in $E$ gilt.

\inputaufgabe
{}
{

Es sei

\mavergleichskette
{\vergleichskette
{ I
}
{ \subseteq }{ \R
}
{  }{
}
{    }{
}
{   }{
}
}
{}{}{}
ein
\definitionsverweis {offenes Intervall}{}{}
und sei
\maabbdisp {p} {E = I \times \R^n } { I
} {}
das triviale Vektorbündel über $I$, das wir mit dem
\definitionsverweis {Standardskalarprodukt}{}{}
als
\definitionsverweis {riemannsches Vektorbündel}{}{}
auffassen. Zeige, dass ein
\definitionsverweis {linearer Zusammenhang}{}{}
$\nabla$ auf $E$ genau dann
\definitionsverweis {metrisch}{}{}
ist, wenn die
\definitionsverweis {Christoffelsymbole}{}{}
die Bedingung
\zusatzklammer {alles bezieht sich auf die einzige Ableitungsrichtung $\partial$} {} {}

\mavergleichskettedisp
{\vergleichskette
{ \Gamma_j^k
}
{ =} { - \Gamma_k^j
}
{ } {
}
{ } {
}
{ } {
}
}
{}{}{}
für

\mavergleichskette
{\vergleichskette
{ 1
}
{ \leq }{ j,k
}
{ \leq }{ n
}
{    }{
}
{   }{
}
}
{}{}{}
erfüllen.

}
{} {}

\inputaufgabe
{}
{

Es sei
\maabb {p} { E } { M
} {}
ein
\definitionsverweis {riemannsches Vektorbündel}{}{}
auf einer
\definitionsverweis {differenzierbaren Mannigfaltigkeit}{}{}
$M$ und sei $\nabla$ ein
\definitionsverweis {linearer Zusammenhang}{}{}
auf $E$. Es sei vorausgesetzt, dass $\nabla$ lokal für die
\definitionsverweis {Standardvektorfelder}{}{}
$\partial_1  , \ldots , \partial_n$ und
\definitionsverweis {Basisschnitte}{}{}

\mathl{s_1  , \ldots , s_r}{} von $E$ die Eigenschaft

\mavergleichskettedisp
{\vergleichskette
{ \partial_i \left\langle  s_j  , s_k   \right\rangle
}
{ =} { \left\langle \nabla_{\partial_i } s_j   ,  s_k   \right\rangle  + \left\langle s_j  ,  \nabla_{\partial_i } s_k    \right\rangle
}
{ } {
}
{ } {
}
{ } {
}
}
{}{}{}
für alle $i,j,k$ erfüllt. Zeige, dass $\nabla$
\definitionsverweis {metrisch}{}{}
ist.

}
{} {}

Für die folgende Aufgabe vergleiche
[[Eingebettete Mannigfaltigkeit/Inverse Karte/Tangentialbündel/Partielle Ableitung/Aufgabe|Aufgabe 10.16]]

\inputaufgabe
{}
{

Es sei

\mavergleichskette
{\vergleichskette
{ Y
}
{ \subseteq }{ \R^n
}
{  }{
}
{    }{
}
{   }{
}
}
{}{}{}
eine
\definitionsverweis {abgeschlossene Untermannigfaltigkeit}{}{}
und
\maabb {\varphi} { V } { U \subseteq Y
} {}
eine $C^2$-diffeomorphe Parametrisierung einer offenen Teilmenge

\mavergleichskette
{\vergleichskette
{ U
}
{ \subseteq }{ Y
}
{  }{
}
{    }{
}
{   }{
}
}
{}{}{}
durch eine offene Teilmenge

\mavergleichskette
{\vergleichskette
{ V
}
{ \subseteq }{ \R^d
}
{  }{
}
{    }{
}
{   }{
}
}
{}{}{,}
wobei wir $\varphi$ auch als Abbildung nach $\R^n$ auffassen. Zeige, dass ein kommutatives Diagramm

\mathdisp {\begin{matrix}  V  & \stackrel{ \partial_i }{\longrightarrow} &  TV & \stackrel{ T( \partial_j) }{\longrightarrow} & TTV & \\ \!\!\!\!\! \,\, \, \varphi \downarrow & & \!\!\!\!\! T(\varphi) \downarrow & &  \downarrow TT(\varphi) \!\!\!\!\! & & \\  U  & \stackrel{  }{\longrightarrow} &  TU & \stackrel{  }{\longrightarrow} & TTU & \!\!\!\!\!  \\ \end{matrix}} {  }

vorliegt. Wie kann man die Abbildungen in den Diagonalen beschreiben?

}
{} {}

  \zwischenueberschrift{Aufgaben zum Abgeben}

\inputaufgabe
{2}
{

Es sei

\mavergleichskette
{\vergleichskette
{ I
}
{ \subseteq }{ \R
}
{  }{
}
{    }{
}
{   }{
}
}
{}{}{}
ein
\definitionsverweis {reelles Intervall}{}{}
und sei

\mavergleichskettedisp
{\vergleichskette
{ g(t)
}
{ =} { 1+t^4
}
{ } {
}
{ } {
}
{ } {
}
}
{}{}{,}
die wir als eine
\definitionsverweis {riemannsche Metrik}{}{}
auf $I$ interpretieren. Bestimme die
\definitionsverweis {vertikale Ableitung}{}{}
\zusatzklammer {zu dem zugehörigen
\definitionsverweis {Levi-Civita-Zusammenhang}{}{}} {} {}

\mathl{\nabla_\partial \left(  3e^{-t^2} + t  \cos \left(  t  \right) - t^5  \right)}{.}

}
{} {}

\inputaufgabe
{2}
{

Berechne $\Gamma^1_{12},  \Gamma^2_{12},  \Gamma^1_{22},  \Gamma^2_{22}$ in
[[Halbebene/Poincaré/Riemannsche Geometrie/Beispiel|Beispiel 18.8]].

}
{} {}

\inputaufgabe
{3}
{

Es sei ${\mathbb H}$ die
\definitionsverweis {Halbebene von Poincaré}{}{}
mit der zugehörigen
\definitionsverweis {riemannschen Struktur}{}{}
und dem zugehörigen
\definitionsverweis {Levi-Civita-Zusammenhang}{}{.}
Bestimme die
\definitionsverweis {vertikale Ableitung}{}{}
\mathl{\nabla_{xy\partial_1 -2y^3 \partial_2 } (G)}{} zum Vektorfeld

\mavergleichskettedisp
{\vergleichskette
{ G(x,y)
}
{ =} { \left(  x^4-2y^2  , \,   3 x^3-  \ln  \left(  1+x^2y^2 \right)                   \right)
}
{ } {
}
{ } {
}
{ } {
}
}
{}{}{.}

}
{} {}

\inputaufgabe
{2}
{

Es sei $\nabla$ ein
\definitionsverweis {torsionsfreier}{}{}
\definitionsverweis {linearer Zusammenhang}{}{}
auf einer offenen Menge

\mavergleichskette
{\vergleichskette
{ U
}
{ \subseteq }{ \R^n
}
{  }{
}
{    }{
}
{   }{
}
}
{}{}{.}
Durch wie viele unabhängige
\definitionsverweis {Christoffelsymbole}{}{}
wird er
\zusatzklammer {in Abhängigkeit von $n$} {} {}
bestimmt?

}
{} {}

 [[Kategorie:Latexseite]]
